\scanpage{12}
si $s'$ est son point fermé, on trouve aussi un homomorphisme de restrictio\unsure{n}
\[
K_{\mathrm{BT}}(S',G)\handwritten{\longrightarrow}K_{\mathrm{BT}}(s',G)\ ,\tag*{\handwritten{(2)}}
\]
en notant que $G$ opère sur $s'$ (de façon pas nécessairement fidèle). Or,$k'=k(s')$ étant supposé \deleted{complet}\inserted{parfait}, on a un foncteur \uline{covariant} naturel
\[
\mathrm{BT}(s')\deleted{\unsure{G)}}\handwritten{\unsure{\mathord{-}|\longrightarrow}}\mathrm{Modlib}(W(k'))\ ,\quad
\handwritten{\text{et \unsure{$=$|même} vers }\mathrm{Cris}(k')}
\]
et comme $G$ opère sur $k'$, donc sur $W(k')$, \deleted{\unsure{on}}le foncteur précédent induit trivialement un foncteur
\[
\mathrm{BT}(s',G)\handwritten{\longrightarrow}\mathrm{Modli}\deleted{\unsure{g}}\mathrm b(W(k'),G),
\]
d'où par suite un homomorphisme
\[
K_{\mathrm{BT}}(s',G)\handwritten{\longrightarrow}K(\deleted{\unsure{Mo}}W(k'),G)\ .\tag*{\handwritten{(3)}}
\]
Composant (1),(2) et (3), on trouve un homomorphisme canonique
\[
K_{\mathrm{BT}}(K;K')\handwritten{\longrightarrow}K(W(k'),G)\ .\tag*{\handwritten{(4)}}
\]
Notons d'ailleurs que si $k$ est algébriquement clos, alors $G$ opère trivialement sur $k'$ donc sur $W(k')$, et on a alors
\[
K(W(k'),G)=R_{W(k')}(G)\ ;
\]
dans le cas général, il y a lieu d'introduire le groupe d'intertie $I\handwritten{\subset}G$, et on trouve un homomorphisme
\[
K(W(k'),G)\handwritten{\longrightarrow}K(W(\deleted{\unsure{j}}k'),I)^{\inserted{\handwritten{G\deleted{\unsure{/I}}}}}\handwritten{\simeq}R_{W(k')}(I)^{\handwritten{G/I}},
\]
\margin{\handwritten{\uline{NB} $G\deleted{\unsure{/I}}$ opère sur $R$ via ses opérations sur $I$ \uline{et} sur $W(k')$ par transport de structure}}
% slanted
\margin{\handwritten{On peut \unsure{distinguer|dire} suivant les \unsure{différentes pentes|différents points} \ldots.}}
d'où grâce à (4)
\[
K_{\mathrm{BT}}(K,K')\handwritten{\longrightarrow}R_{W(k')}(I)^{\handwritten{G\deleted{\unsure{/I}}}}
\handwritten{\big[}\handwritten{\longrightarrow}\mathrm{Cent}_{\handwritten{G}}(I,W(\handwritten{k}'))\ ,\tag*{\handwritten{(5)}}
\]
la dernière flèche étant la flèche trace. L'homomorphisme composé est appelé encore l'homomorphisme trace, l'image de $\mathrm{c}\handwritten{\ell}(M)$ dahs le dernier membre est appelé le \uline{caractère} du groupe de B.T. $M$ . On voit facilement que dans un sens évident, il ne dépend pas du choix du corps de bonne réduction virtuelle $K'$: la fonction caractère est alors envisagée comme fonction sur le \inserted{sous-}groupe d'inertie (infini) de $\mathrm{Gal}(\overline K/K)$.

Supposons que $M_{K'}$ soit \uline{constant}, de sorte que \deleted{\unsure{$M$ est un f $M$}} $T_p(M)$ est un faisceau $p$-adique localement constant (au sens technique
% archive sheet number: 11

